\section{Introduction}
Trong bối cảnh điện khí hóa đang trở thành khuynh hướng tất yếu trong ngành giao thông vận tải và năng lượng tái tạo, nhu cầu về các dòng máy điện có mật độ công suất và hiệu suất vượt trội ngày càng trở nên cấp thiết. Trong đó, động cơ nam châm vĩnh cửu hướng trục (AFPM) đã khẳng định vị thế là một trong những công nghệ then chốt nhờ những ưu thế đặc biệt về cấu trúc \cite{jenkins2023innovations,gieras2008axial,capponi2012recent,xu2021comparative}. 

Tuy nhiên, khác biệt hoàn toàn với các dòng máy điện hướng kính truyền thống, AFPM sở hữu đặc thù hình học với sự phân bổ từ trường biến thiên phi tuyến và phức tạp trong không gian ba chiều \cite{vatani2025axial,ni2013investigation}. Đặc tính này khiến các mô hình 2D đơn giản hóa không còn đủ khả năng đảm bảo độ tin cậy trong việc phân tích và đánh giá chính xác các đặc tính vận hành. 

Để đạt được độ chính xác cần thiết trong thực tế, các nhà thiết kế bắt buộc phải sử dụng đến phương pháp Phần tử hữu hạn 3D (3D-FEM). Mặc dù vậy, thách thức nảy sinh khi các quy trình thiết kế hiện đại thường đi kèm với bài toán tối ưu hóa đa mục tiêu sử dụng các thuật toán tìm kiếm toàn cục \cite{huner2023design,subotic2019weight,liu2025optimization}. Những thuật toán này đòi hỏi hàng ngàn chu kỳ đánh giá mô hình lặp lại, biến chi phí tài nguyên và thời gian thực thi cực lớn của 3D-FEM truyền thống trở thành một nút thắt cổ chai kỹ thuật, cản trở tiến độ nghiên cứu và phát triển sản phẩm.

Các phương pháp tham số tập trung (Lumped Parameter Method - LPM), hay còn được biết đến với tên gọi mạng từ trở tương đương (Magnetic Equivalent Circuit - MEC), đã được nghiên cứu và phát triển từ rất sớm như một hướng tiếp cận kinh điển nhằm tối ưu hóa tốc độ tính toán thông qua việc mô hình hóa không gian bằng một hệ thống mạch từ cấu tạo từ các từ trở và nguồn sức từ động. Nhờ ưu thế vượt trội về mặt thời gian thực thi, phương pháp này đã được ứng dụng thành công và rộng rãi trên nhiều dòng máy điện từ truyền thống như động cơ nam châm vĩnh cửu bề mặt (SPMSM) \cite{manju2018magnetic,cui2025equivalent,zhang2019equivalent,yang2011surface}, nam châm chìm (IPMSM) \cite{kemmetmuller2014modeling,zhu2009analytical,tariq2010iron} cho đến các cấu trúc AFPM hiện đại \cite{zhang20213,6971543,6600909,7869440}. Tuy nhiên, việc xây dựng mô hình trong các phương pháp này thường bị ràng buộc bởi các giả thiết hình học và quy ước dòng từ dựa trên kinh nghiệm, khiến các quy luật phân bố từ trường nằm ngoài phạm vi dự tính dễ bị bỏ sót, gây ra những hạn chế đáng kể khi đánh giá các cấu trúc có trường từ phức tạp hoặc đòi hỏi độ chính xác cục bộ cao.

Một nhánh nghiên cứu gần đây dựa trên việc rời rạc hóa không gian tính toán thành duy nhất một loại phần tử hình học thông qua lưới phi tham số (non-parametric mesh) đã mở ra hướng đi mới cho kỹ thuật mạch từ tương đương thay thế (MEC). Phương pháp này được giới thiệu ban đầu thông qua việc sử dụng các phần tử hình chữ nhật trong hệ tọa độ Cartesian để tự động hóa quá trình tạo lưới \cite{9473194}. Các công trình tiêu biểu tiếp theo đã phát triển kỹ thuật này bằng cách sử dụng các phần tử hình tam giác \cite{10018857,ouagued2016comparison}, hình quạt \cite{van2026development,lu2023hybrid,10599821} cho phép mô hình hóa các dòng máy điện hướng kính một cách hiệu quả. Tuy nhiên, các nghiên cứu này chỉ dừng lại trên mặt phẳng tính toán hai chiều (2D), do đó việc khai thác kĩ thuật này trong không gian 3 chiều còn nhiều tiềm năng.

Xuất phát từ bối cảnh đó, bài báo này đề xuất một phương pháp mạng từ trở sinh lưới trong không gian 3 chiều (3D-MBGRN) dựa trên hệ tọa độ trụ cho các máy điện quay đối xứng. Phương pháp thực hiện rời rạc hóa miền tính toán dựa trên cấu trúc lưới đồng nhất để chia không gian thành các phần tử hình lục diện, cho phép mô tả đầy đủ sự phân bố từ trường cục bộ tại mọi vị trí. Để nâng cao hiệu năng tính toán, kỹ thuật dịch chỉ số (index-shifting) được kết hợp cùng các điều kiện biên chu kỳ và phản chu kỳ nhằm thu gọn miền tính toán và mô phỏng chuyển động quay của rotor mà không làm thay đổi cấu trúc lưới cơ sở, loại bỏ hoàn toàn chi phí tái tạo lưới hoặc nội suy ranh giới trượt. Bên cạnh đó, bộ giải lặp điểm cố định tích hợp thuật toán suy giảm thích nghi hệ số thư giãn vật liệu (Adaptive Material Decay) kết hợp cơ chế quay lui nghiệm (backtracking) giúp triệt tiêu rủi ro rơi vào bẫy trì trệ hội tụ và dao động số học.

Các đóng góp chính của bài báo bao gồm:
\begin{enumerate}
    \item Đề xuất khung mô hình mạng từ trở 3D sinh lưới lục diện (3D-MBGRN) trên hệ tọa độ trụ kết hợp thuật toán gán vật liệu tự động theo vị trí (point-in-region), áp dụng cho các dòng máy điện quay đối xứng.
    \item Ứng dụng kỹ thuật dịch chỉ số lưới (index-shifting) nhằm mô phỏng chuyển động quay của rotor trên nền lưới tĩnh, nâng cao tốc độ thực thi transient nhờ loại bỏ hoàn toàn các tác vụ tái lập ma trận hình học và phép nội suy mặt trượt phức tạp.
    \item Xây dựng bộ giải lặp phi tuyến thích nghi với thuật toán Adaptive Material Decay kết hợp cơ chế quay lui (rollback), đảm bảo hội tụ số học bền vững và giảm thiểu hiện tượng bẫy giả hội tụ (stagnation trap) phổ biến trong các bộ giải thư giãn cố định.
    \item Thực hiện kiểm chứng toàn diện mô hình đề xuất với 3D-FEM trên các trạng thái vận hành không tải và có tải, chứng minh độ hội tụ đơn điệu theo mật độ ô lưới và khả năng tiệm cận độ chính xác của 3D-FEM trên toàn bộ các đặc tính transient (kể cả cogging torque), đồng thời giảm kích thước ma trận hệ thống (DoF) $1.4~\text{lần}$ và cải thiện tốc độ tính toán từ $6$ đến $20$ lần ở quy mô lưới tương đương.
\end{enumerate}

Bài báo được trình bày thành 7 phần: Phần I giới thiệu tổng quan bối cảnh và mục tiêu nghiên cứu; Phần II trình bày mô hình đề xuất 3D-MBGRN bao gồm quy trình rời rạc hóa miền tính toán, cấu trúc phần tử lục diện, các điều kiện biên và kỹ thuật mô phỏng chuyển động quay; Phần III mô tả việc triển khai số trị với các phương trình bảo toàn cục bộ, kỹ thuật lắp ghép ma trận hệ thống và bộ giải phi tuyến thích nghi; Phần IV trình bày các công thức hậu xử lý để trích xuất các đại lượng đặc tính điện từ; Phần V thực hiện kiểm chứng mô hình đề xuất thông qua việc đối chiếu toàn diện với 3D-FEM ở cả hai trạng thái không tải và có tải định mức; Phần VI thảo luận chuyên sâu về đặc tính hội tụ của bộ giải và phân tích độ nhạy lưới NRMSE; và Phần VII đưa ra các kết luận chính cùng định hướng phát triển của nghiên cứu